\documentclass[a4paper,10pt,oneside,fleqn]{jsbook}
%
\usepackage{amsmath,amssymb,bm}
\usepackage{bm}
\usepackage{graphicx}
\usepackage{ascmac}
\usepackage{makeidx}
\usepackage{txfonts}
\usepackage{indentfirst}
\usepackage{indent}
\usepackage{booktabs}
\usepackage{comment}
\usepackage{cite}
\usepackage{subfigure}
\usepackage{alltt}
\usepackage{eclbkbox,fancybox}
\usepackage{float}
%
\newcounter{program}
\newenvironment{program}%
{\vspace{0.5\baselineskip}\VerbatimEnvironment%
%\begin{breakbox}\setlength{\baselineskip}{0pt}\begin{Verbatim}}%
\begin{breakbox}\setlength{\baselineskip}{.25\normalbaselineskip}\begin{Verbatim}}%
{\end{Verbatim}\end{breakbox}\vspace{0.8\baselineskip}}

\makeindex
%
\newcommand{\diff}{\mathrm{d}}  %微分記号
\newcommand{\divergence}{\mathrm{div}\,}  %ダイバージェンス
\newcommand{\grad}{\mathrm{grad}\,}  %グラディエント
\newcommand{\rot}{\mathrm{rot}\,}  %ローテーション
%
\setlength{\textwidth}{\fullwidth}
\setlength{\textheight}{44\baselineskip}
\addtolength{\textheight}{\topskip}
\setlength{\voffset}{-0.6in}
\setlength{\mathindent}{1.5cm} %数式のインデント設定
\setcounter{secnumdepth}{2} %subsubsectionまで番号付け
\setcounter{tocdepth}{4} %目次の深さ設定
\renewcommand\citemid{; } %citeのオプション

%subfigure
\renewcommand{\subfigtopskip}{5pt}	% 図の上の隙間。上図の副題と下図の間
\renewcommand{\subfigbottomskip}{5pt} % 図の下の隙間。副題と本題の間
\renewcommand{\subfigcapskip}{10pt}	% 図と副題の間
\renewcommand{\subcapsize}{\small} % 副題の文字の大きさ

%
\usepackage{atbegshi}
\AtBeginShipoutFirst{\special{pdf:tounicode EUC-UCS2}}
\usepackage[dvipdfm,bookmarks=true,bookmarksnumbered=true]{hyperref}
%
\begin{document}

\begin{titlepage}
\begin{center}
\vspace*{2cm}
{\huge \textbf{Performance Monitor Library}}\\
\vspace{1cm}

{\large \textbf{Ver. 1.0}}\\
\vspace{1.0cm}

{\large \textbf{Advanced Visualization Research Team}\\
\large \textbf{Advanced Institute of Computational Science}\\
\large \textbf{RIKEN}\\
\vspace{1cm}
}

{7-1-26, Minatojima-minami-machi, Chuo-ku, Kobe, Hyogo, 650-0047, Japan}\\
\vspace{0.5cm}

\url{http://aics.riken.jp/}\\

\vspace{1.0cm}




April 2012\\
\vspace{3cm}

\includegraphics[width=4cm,bb=-80 0 220 500]{RIKEN_logo_300x500.eps}

\end{center}
\end{titlepage}
\newpage

%
\frontmatter

\begin{tabular}{llllr}
First Edition  &  Version 1.0  & 28 Apr. & 2012



\end{tabular}

\vspace{10cm}

\begin{center}
\includegraphics[width=4cm]{kei.eps}
\end{center}

\vspace{2cm}

\begin{description}
\item[ ] \textbf{COPYRIGHT}\\
(c) Copyright RIKEN 2012. All rights reserved.\\

%\item[ ] \textbf{DISCLAIMER}\\
%You shall comply with the conditions of the license when you use this program.\\
%The license is available at http://vcad-hpsv.riken.jp/permission.html
\end{description}
%

\tableofcontents
%
%
\mainmatter

%
\chapter{計算性能モニタクラス}
{\begin{abstract}
本章では，実行時の計算性能を測定するモジュール機能について説明する．
\end{abstract}
\include{PM_Monitor}

%
\chapter{アップデート情報}
\begin{abstract}
本ユーザガイドのアップデート情報について記す．
\end{abstract}

{\small

%

%
\subparagraph{Version 1.1\hspace{1cm}2012/5/2}

\begin{description}
\begin{indentation}{3zw}{0zw}
\item[-] Version 1.1 リリース．
\item[-] ヘッダの修正
\end{indentation}
\end{description}
\vspace{3mm}


\subparagraph{Version 1.0\hspace{1cm}2012/4/28}

\begin{description}
\begin{indentation}{3zw}{0zw}
\item[-] Version 1.0 リリース．
\item[-] 利用方法のメモを追加 How\_to\_use\_PMlib.pdf
\end{indentation}
\end{description}
\vspace{3mm}

} % end of small

%
%\bibliographystyle{junsrt}
%\bibliography{PMlib}

\newpage
\printindex
%
%
\end{document}